% ---------------------------------------------
% 第3章：新電力の倒産構造と仮説（第1章との整合版）
% ---------------------------------------------

\chapter{新電力の倒産構造と仮説}
\label{chap:hypothesis}

本章では、第1章で整理した新電力事業者が直面する五つのリスク
（価格変動リスク、インバランスリスク、需給逼迫リスク、
制度変更リスク、資本制約リスク）を踏まえ、
これらがどのように倒産リスクへと結びつくのかを理論的に整理する。
さらに、この倒産構造に基づき、本研究で検証する仮説を提示する。

本章で提示する倒産構造は、第1章で整理した制度的背景と、第2章で示した問題意識・研究目的を統合したものであり、新電力企業の倒産リスクがどのような構造原理によって決定されるのかを理論的に再構成する役割を持つ。さらに、本章で整理する四要素モデルは、日本市場だけでなく制度構造の異なる英国市場にも適用可能であり、倒産リスクの構造原理が国境を越えて普遍的に成立するかを検証するための基礎となる。

本章では、以上の整理を踏まえ、後節で四つの仮説（H1〜H4）を提示する。

第1章で示したように、新電力の倒産は
\begin{itemize}
    \item 外部ショック（市場価格・制度コスト）が大きく、
    \item 財務指標が良好でも短期間で倒産に至り、
    \item 親会社・自治体などの支援能力の有無が決定的であり、
    \item 制度コストの急増により倒産閾値が変動する
\end{itemize}
という特徴を持つ。

これらの特徴は、従来の財務比率モデルでは説明しきれず、
新電力市場に固有の構造的要因を踏まえた理論的整理が必要となる。

% ---------------------------------------------
% 3.1 新電力倒産の特徴：一般倒産研究との相違点
% ---------------------------------------------

\section{新電力倒産の特徴：一般倒産研究との相違点}

本節では、新電力企業の倒産事例が一般企業の倒産研究で想定される
「財務指標の悪化が倒産を規定する」という通説と大きく異なる構造を示すことを整理する。
この事実整理は、後節で提示する仮説（H1〜H4）が自明ではなく、
新電力市場に特有の倒産構造を明らかにするための基礎となる。

\subsection*{(1) 財務指標が健全でも倒産する外部ショック型倒産}

一般企業の倒産研究では、倒産は財務指標の悪化（利益率低下、債務超過、流動性不足）
によって徐々に進行するとされる。
しかし新電力市場では、財務指標が健全であっても外部ショックにより
短期間で倒産に至る事例が多数報告されている。

日本市場では、2021年のJEPX価格高騰とインバランス料金急増により、
直前期まで黒字であった新電力A社が数週間で資金繰り破綻に至った
\parencite{TDB2021PowerCrisis,TSR2021PowerCrisis}。
同様のショックを受けたにもかかわらず、自治体支援を受けたB社は倒産を回避し、
支援主体が存在しなかったC社は撤退に追い込まれた
\parencite{DenkiShimbun2021Crisis,CityCouncil2021Support}。

英国市場でも同様の構造が観察される。
燃料価格高騰期において、Bulb Energy は価格キャップ制度により
調達費用の上昇を料金へ転嫁できず倒産した一方、
Octopus Energy は親会社の資本支援と厚いヘッジ契約により同じショックを乗り切った
\parencite{Ofgem2021BulbSLR,Ofgem2021OctopusTakeover,UKGov2021BulbSupport}。

これらの事例は、倒産が財務指標の線形的悪化ではなく、
外部ショックへの曝露度（$\sigma$）によって決定される場合があることを示している。

\subsection*{(2) 支援能力の有無が倒産回避の決定要因となる構造}

一般企業の倒産研究では、親会社・自治体・金融機関の支援能力は補助的要因とされ、
倒産確率を規定する主要因として扱われることは少ない。
しかし新電力市場では、支援能力（$S$）の有無が倒産回避の可否を左右する
決定的要因となっている。

日本市場では、自治体や親会社による短期資金支援が倒産回避の分岐点となった事例が複数報告されている。
英国市場でも、Octopus Energy が親会社の資本支援により倒産を回避した一方、
支援主体を持たなかった小売事業者は同じショックで倒産に至った。

この構造は、支援能力が倒産閾値（$D$）を実質的に引き下げる外生的緩和要因として機能することを示している。

\subsection*{(3) 制度変更により倒産閾値が短期間で急増する非線形構造}

一般企業では、負債返済額やキャッシュアウトは比較的安定しており、
倒産閾値（$D$）は急激に変動しない。
しかし新電力市場では、制度変更や需給逼迫により倒産閾値が短期間で急増する。

日本市場では、需給調整金やインバランス負担が制度変更により急増し、
平常時とは異なる倒産ラインが形成されたと指摘されている。
英国市場でも、価格キャップ制度の改定や SLR 発動時の追加負担が
倒産閾値を押し上げ、複数の小売事業者が短期間で倒産に至った。

このような制度依存的な倒産閾値の急増は、
一般企業には見られない新電力特有の非線形構造である。

\subsection*{(4) 複数リスクが同時悪化した瞬間に倒産が顕在化する急峻反応}

一般企業の倒産研究では、倒産確率は財務指標の悪化に伴い徐々に上昇する。
しかし新電力市場では、複数のリスクが同時に悪化した瞬間に
倒産が急峻に顕在化する「急峻反応」が観察される。

日本市場では、価格高騰・インバランス急増・需給逼迫・ヘッジ不足・資本制約が
同時に悪化した瞬間に倒産が発生した事例が多数報告されている。
英国市場でも、燃料価格高騰と制度負担増加が同時に発生した際に
複数の小売事業者が連鎖的に倒産した。

この急峻反応は、収益性（$\mu$）、外部ショック感応度（$\sigma$）、
支援能力（$S$）、倒産閾値（$D$）の multiplicative な相互作用によって生じる構造であり、
一般倒産研究の線形的悪化モデルでは説明できない。

\subsection*{(5) 本節のまとめ}

以上の事実整理より、新電力企業の倒産は
\begin{itemize}
    \item 財務指標が健全でも倒産する（H1 の根拠）
    \item 支援能力の有無が倒産回避の決定要因となる（H2 の根拠）
    \item 制度変更により倒産閾値が急増する（H3 の根拠）
    \item 複数リスク同時悪化により倒産が急峻に顕在化する（H4 の根拠）
\end{itemize}
という一般倒産研究とは異なる構造を示す。

したがって、次節で提示する仮説（H1〜H4）は自明な命題ではなく、
新電力市場に特有の倒産構造を明らかにするための検証価値の高い仮説である。



% ---------------------------------------------
% 3.2 第1章の五つのリスクの統合：四要素モデルへの抽象化
% ---------------------------------------------

\section{五つのリスクの統合：四要素モデルへの抽象化}

第1章で整理した五つのリスクは、倒産構造の観点から整理すると、
以下の四つの要素に抽象化できる。

\begin{enumerate}
    \item \textbf{収益性 $\mu$}  
          （価格変動・需給逼迫による収益の変動性）
    \item \textbf{外部ショック感応度 $\sigma$}  
          （スポット依存度・インバランス曝露度）
    \item \textbf{支援能力 $S$}  
          （親会社・自治体・金融機関の支援の有無）
    \item \textbf{倒産閾値 $D$}  
          （制度変更・需給逼迫時のキャッシュアウト増加）
\end{enumerate}

第1章で整理した五つのリスクは、それぞれが独立して倒産リスクに影響するのではなく、
新電力企業の経営構造において相互に関連しながら作用する。

この相互作用的な倒産構造は、第1章および第2章で示した日本・英国の倒産事例とも整合的である。
日本市場では、2021年のJEPX価格急騰により、直前期まで黒字であった新電力A社が
数週間で資金繰り破綻\footnote{帝国データバンクおよび東京商工リサーチは、2021年以降の新電力倒産の特徴として「財務指標が健全でも外部ショックにより短期間で資金繰りが破綻する」外部ショック型倒産の存在を指摘している。}
に至った\parencite{TDB2021PowerCrisis,TSR2021PowerCrisis}。
同じショックを受けたにもかかわらず、自治体支援を受けたB社は倒産を回避し、
支援主体が存在しなかったC社は撤退\footnote{電気新聞および自治体議会資料によれば、自治体との包括契約に基づく短期資金支援の有無が倒産回避の可否を左右したと報告されている。}
に追い込まれた\parencite{DenkiShimbun2021Crisis,CityCouncil2021Support}。

英国市場でも同様の構造が観察される。
燃料価格高騰期において、Bulb Energy は価格キャップ制度により調達費用の上昇を料金へ転嫁できず倒産\footnote{新電力の倒産事例では、親会社・自治体・金融機関などの支援主体の存在が倒産回避の可否を決定づけると複数の倒産レポートで指摘されている。}
した一方、
Octopus Energy は親会社の資本支援と厚いヘッジ契約により同じショックを乗り切った\footnote{Ofgem の公式記録では、価格キャップ制度の改定が小売事業者のキャッシュフローを圧迫し、倒産ラインを大きく押し上げた事例が複数報告されている。}
\parencite{Ofgem2021BulbSLR,Ofgem2021OctopusTakeover,UKGov2021BulbSupport}。

これらの事例は、倒産が財務指標の線形的悪化ではなく、
収益性（$\mu$）、外部ショック感応度（$\sigma$）、支援能力（$S$）、倒産閾値（$D$）という
四つの構造的要因の相互作用によって決定されることを示している。
したがって、五つのリスクを四要素モデルへと抽象化することは、
新電力企業の倒産構造を理論的に整理するうえで不可欠である。

これらのリスクを倒産構造の観点から整理すると、四つの構造的要素
（収益性 $\mu$、外部ショック感応度 $\sigma$、支援能力 $S$、倒産閾値 $D$）
へと抽象化できる。

表 \ref{tab:risk_mapping} は、第1章で示した五つのリスクと、
本章で導入する四要素モデル（$\mu, \sigma, S, D$）の対応関係を整理したものである。
この対応関係により、実態として観察されるリスク（第1章）と、
理論モデルで扱う変数（第4章）がどのように結びつくかが明確になる。

\begin{table}[H]
\centering
\caption{第1章の五つのリスクと四要素モデル（$\mu, \sigma, S, D$）の対応関係}
\label{tab:risk_mapping}
\begin{tabular}{p{4cm} p{4cm} p{6cm}}
\toprule
\textbf{第1章のリスク} & \textbf{対応する要素} & \textbf{倒産構造における意味} \\
\midrule
価格変動リスク & 収益性 $\mu$・外部ショック感応度 $\sigma$ &
収益の変動性と市場価格への曝露度を規定する。 \\

インバランスリスク & 外部ショック感応度 $\sigma$・倒産閾値 $D$ &
制度コストの急増がショック時の損失を拡大し、倒産ラインを押し上げる。 \\

需給逼迫リスク & 収益性 $\mu$・倒産閾値 $D$ &
需給逼迫時の調達コスト増加がキャッシュアウトを増大させる。 \\

制度変更リスク & 倒産閾値 $D$ &
制度的キャッシュアウトの急増により倒産ラインが非線形に変動する。 \\

資本制約リスク & 支援能力 $S$ &
親会社・自治体・金融機関の支援の有無が倒産回避能力を決定する。 \\
\bottomrule
\end{tabular}
\end{table}

この対応関係から明らかなように、四要素モデルは第1章で示した五つのリスクを
単に圧縮したものではなく、倒産に直接影響する構造的因子として再整理したものである。
特に、外部ショック感応度 $\sigma$ と倒産閾値 $D$ は、
市場環境の急変により大きく変動し、財務指標では捉えにくい非線形性を生む点で重要である。

以下では、この四要素モデルを基礎として、
新電力企業の倒産がどのようなメカニズムで発生するのかを理論的に整理し、
multiplicative な倒産構造へと議論を進める。


この四要素は、第1章の五つのリスクを
「倒産に直接影響する構造的因子」として再整理したものである。

特に、$\sigma$ と $D$ は市場環境により急変し、
財務指標では捉えにくい非線形性を生む点で重要である。

なお、四要素モデル（$\mu, \sigma, S, D$）は新電力市場に特有の制度的特徴を踏まえて構築したものであるが、その構造は日本市場に限定されるものではない。英国市場でも、価格キャップ制度の変更や SLR 発動時の負担増加が倒産閾値 $D$ を押し上げ、外部ショック感応度 $\sigma$ が企業価値のボラティリティに直結する構造が確認されている。したがって、四要素モデルは制度差を超えて適用可能な普遍的枠組みであり、後章の実証分析において日英両市場で検証する意義がある。

% ---------------------------------------------
% 3.3 multiplicative 倒産構造：五つのリスクが同時に悪化する現象
% ---------------------------------------------

\section{multiplicative 倒産構造：五つのリスクが同時に悪化する現象}

新電力企業の倒産事例（TDB・TSR）は、
以下のような「複数リスクの同時悪化」が倒産の引き金となっていることを示す。

\begin{itemize}
    \item JEPX 価格高騰（外部ショック）
    \item インバランス料金急増（制度コスト）
    \item 需給逼迫（需給リスク）
    \item ヘッジ不足（事業運営リスク）
    \item 資本制約（財務リスク）
\end{itemize}

これらは単独では倒産に至らないが、
同時に悪化すると短期間で資金繰りが破綻する。

この構造は、Merton（1974）や Shumway（2001）が示した
「multiplicative → log → linear」という倒産構造と整合的である。

したがって、新電力の倒産構造は、
第1章の五つのリスクが multiplicative に作用する現象として理解できる。

この multiplicative 構造は、制度的キャッシュアウトが急増する新電力市場において特に顕著である。価格変動、インバランス、需給逼迫、制度変更、資本制約といった複数のリスクが同時に悪化すると、倒産閾値 $D$ が短期間で跳ね上がり、外部ショック感応度 $\sigma$ が企業価値のボラティリティを急増させる。こうした制度依存的な multiplicative 構造は、日本市場だけでなく英国市場でも観察されており、倒産リスクの構造原理が国や制度を超えて共通する可能性を示唆する。

% ---------------------------------------------
% 3.4 新電力市場における倒産構造：四要素モデルの位置づけ
% ---------------------------------------------


\section{新電力市場における倒産構造：四要素モデルの位置づけ}

以上の議論を踏まえ、本研究では新電力企業の倒産リスクが
以下の四要素の multiplicative な相互作用により決定されると整理する。

\begin{itemize}
    \item 収益性 $\mu$：外部環境に応じて変動するキャッシュフロー期待値
    \item 外部ショック感応度 $\sigma$：市場価格変動への曝露度
    \item 支援能力 $S$：外生的な緩和要因（親会社・自治体・金融機関）
    \item 倒産閾値 $D$：制度変更や需給逼迫により変動する倒産ライン
\end{itemize}

この四要素は、第4章で構築する理論モデル（潜在倒産距離 $Z_{\text{lin}}$）の
構成要素としても自然に対応する。

新電力企業の倒産閾値 $D$ は、制度変更や需給逼迫により短期間で急増するため、DD の解釈には制度的キャッシュアウトの変動を組み込む必要がある。この制度依存的な倒産閾値の変動は英国市場でも同様に観察されており、価格キャップ制度の改定や SLR 発動時の負担増加が倒産距離を急低下させる事例が多数報告されている\footnote{2021年以降の制度変更により、需給調整金やインバランス負担が急増し、平常時とは異なる倒産閾値が形成されたと指摘されている。}
。したがって、DD を新電力企業に適用する際には、制度的ショックが倒産距離に与える影響を明示的に考慮する必要がある。


% ---------------------------------------------
% 3.5 本研究の仮説（H1〜H4）
% ---------------------------------------------

\section{本研究の仮説（H1〜H4）}

\subsection*{本研究の仮説設定に関する位置づけ}

以下で提示する仮説（H1〜H4）は、一見すると
一般的な倒産研究で扱われる命題と類似しているように見える。
しかし、新電力市場の倒産事例や制度的特徴を踏まえると、
これらの仮説は決して自明ではなく、むしろ新電力特有の倒産構造を
明らかにするうえで重要な論点となる。


第一に、新電力企業では財務指標が良好であっても、
外部ショック（JEPX 価格高騰やインバランス料金急増）により
数週間で倒産に至る事例が多数報告されている。
これは「財務指標が倒産を説明する」という一般的前提が成立しないことを意味する。

第二に、親会社・自治体・金融機関の支援能力（$S$）が
倒産回避の決定的要因となる点も、新電力特有の構造である。
一般企業の倒産研究では、支援能力は補助的要因にすぎず、
倒産確率を規定する主要因として扱われることはほとんどない。

第三に、倒産閾値（$D$）が制度変更や需給逼迫により
短期間で急増する点も一般企業には見られない特徴である。
制度的キャッシュアウトが倒産ラインを非線形に押し上げる現象\footnote{英国の SLR（Supplier of Last Resort）制度は、倒産した小売事業者の顧客を他社が引き継ぐ際に追加負担が発生する制度であり、燃料価格高騰期にはこの負担が急増し、倒産ラインを押し上げる要因となったと指摘されている。}
は、
新電力市場に固有のものである。

以上より、本研究の仮説（H1〜H4）は
「一般的な倒産研究の再確認」ではなく、
新電力市場に特有の倒産構造を明らかにするための
非自明かつ検証価値の高い仮説である。


以上の倒産構造を踏まえ、本研究では以下の仮説を設定する。
各仮説については、導出の根拠、理論的背景、および仮説が成立した場合の含意を併せて示す。

\subsection*{仮説1：外部ショック主因仮説（H1）}

\textbf{H1：新電力企業の倒産は、財務指標よりも外部ショック感応度（$\sigma$）によって強く規定される。}

\textit{（導出の根拠）}  
第1章で示したように、2021 年以降の JEPX 価格高騰やインバランス料金の急増は、
黒字決算の企業を数か月で倒産に追い込んだ。
TDB・TSR の倒産事例でも、財務指標が良好であっても外部ショックにより
短期間で資金繰りが悪化した点が繰り返し指摘されている。

\textit{（理論的背景）}  
外部ショックへの曝露度は、Merton 型モデルにおけるボラティリティ $\sigma_V$ に対応し、
multiplicative に倒産確率を押し上げる要因である。

\textit{（含意）}  
H1 が成立すれば、従来の財務比率モデルでは新電力の倒産を説明できず、
外部ショック感応度を組み込んだ倒産距離モデルが必要となる。

\bigskip

\subsection*{仮説2：支援能力優位仮説（H2）}

\textbf{H2：支援能力（$S = 1$）が存在する企業は、財務指標にかかわらず倒産を回避しやすい。}

\textit{（導出の根拠）}  
TSR の倒産事例では、親会社・自治体・金融機関の支援の有無が
倒産回避の分岐点となっていることが明確に示されている。
支援主体が存在する企業は、外部ショック時に追加資金を確保できる。

\textit{（理論的背景）}  
支援能力 $S$ は、構造モデルにおける負債返済能力の外生的緩和要因として機能し、
倒産閾値 $D$ を実質的に引き下げる役割を持つ。

\textit{（含意）}  
H2 が成立すれば、倒産リスクは財務指標だけでは説明できず、
企業のネットワーク構造（親会社・自治体・金融機関）を組み込む必要がある。

\bigskip

\subsection*{仮説3：倒産閾値非線形仮説（H3）}

\textbf{H3：倒産閾値（$D$）は外部ショック時に急増し、倒産確率を非線形的に押し上げる。}

\textit{（導出の根拠）}  
需給調整金やインバランス料金は、制度変更や需給逼迫により
短期間で数倍に跳ね上がることがある。
TDB は「制度的キャッシュアウトの急増が倒産ラインを押し上げた」と指摘している。

\textit{（理論的背景）}  
倒産閾値 $D$ は、Merton 型モデルにおける負債返済額に対応し、
ショック時に急増することで倒産確率を非線形的に高める。

\textit{（含意）}  
H3 が成立すれば、倒産モデルは平常時とショック時で異なるパラメータを持つ必要があり、
線形モデルでは倒産確率の急増を捉えられない。

\bigskip

\subsection*{仮説4：潜在倒産距離の急峻反応仮説（H4）}

\textbf{H4：外部ショック感応度の高さ（H1）、支援能力の欠如（H2）、倒産閾値の急増（H3）が同時に作用する場合、潜在倒産距離 $Z_{\text{lin}}$ のわずかな悪化が倒産確率の急激な上昇につながる。}

\textit{（導出の根拠）}  
新電力の倒産事例では、複数のリスクが同時に悪化した瞬間に
倒産が一気に顕在化する\footnote{英国市場では、価格キャップ制度の改定や SLR 発動時の負担増加が倒産距離を急低下させる事例が多数報告されており、制度的ショックが倒産距離に与える影響を明示的に考慮する必要がある。}
「急峻反応」が観察されている。

\textit{（理論的背景）}  
H1〜H3 の multiplicative な構造が重層的に作用すると、
$Z_{\text{lin}}$ の小さな変化が倒産確率の大きな変化につながる。

\textit{（含意）}  
H4 が成立すれば、新電力の倒産リスクは「連続的に悪化する」のではなく、
「閾値を超えた瞬間に急激に悪化する」ことが示される。
これは第4章で構築する潜在倒産距離モデルの非線形性を支持する。

% ---------------------------------------------
% 3.6 本章のまとめと第4章への接続
% ---------------------------------------------

\section{本章のまとめと第4章への接続}

本章で提示した四要素モデルと仮説（H1〜H4）は、
第4章の理論モデル（潜在倒産距離 $Z_{\text{lin}}$）において形式化される。
特に、$\sigma$ と $D$ の変動性は、
新電力企業の倒産リスクを従来モデルとは異なる形で規定する要因として重要である。

以上の理論的整理により、本研究で構築する潜在倒産距離モデル（$Z_{\text{lin}}$）は、四要素モデル（$\mu, \sigma, S, D$）の構造を形式化し、multiplicative な倒産構造を定量的に評価する枠組みとして位置づけられる。第4章では、この理論モデルを日本企業および英国企業の財務データに適用し、倒産リスクの構造原理が制度差を超えて普遍的に成立するかを検証する。

以上の理論的整理により、本研究で構築する潜在倒産距離モデル（$Z_{\text{lin}}$）は、四要素モデル（$\mu, \sigma, S, D$）の構造を形式化し、multiplicative な倒産構造を定量的に評価する枠組みとして位置づけられる。第4章では、この理論モデルを日本企業および英国企業の財務データに適用し、倒産リスクの構造原理が制度差を超えて普遍的に成立するかを検証する。
